% ============================================================
%  Chapter 9 — Linear Dependence and Independence
% ============================================================
\section{Linear Dependence and Independence}
\label{sec:lindep}

One of the central questions in linear algebra is: can a vector be expressed
as a linear combination of others?  The concepts of \emph{linear dependence}
and \emph{linear independence} give a precise vocabulary for answering this
question, and they underlie the notion of a basis.

% ---- Definitions -------------------------------------------
\begin{defbox}[Definition --- Linear Dependence and Independence]
\label{def:lin-dep-indep}
Let $V$ be a vector space over $\F$ and $S\subseteq V$.

\begin{itemize}
    \item $S$ is \textbf{linearly dependent} (l.d.) if there exist
          \emph{distinct} vectors $v_1,\ldots,v_n\in S$ and scalars
          $\alpha_1,\ldots,\alpha_n\in\F$, not all zero, such that
          \[
              \alpha_1 v_1 + \alpha_2 v_2 + \cdots + \alpha_n v_n = \mathbf{0}.
          \]
    \item $S$ is \textbf{linearly independent} (l.i.) if it is not linearly
          dependent.
\end{itemize}
\end{defbox}

% ---- Finite-set formulation --------------------------------
\begin{tipbox}[Remark --- Linear Independence for Finite Sets]
For a finite set $S=\{v_1,\ldots,v_n\}$, linear independence means:
\[
    \alpha_1 v_1 + \cdots + \alpha_n v_n = \mathbf{0}
    \implies \alpha_1 = \cdots = \alpha_n = 0.
\]
That is, the \emph{only} linear combination equal to $\mathbf{0}$ is the
trivial one.
\end{tipbox}

% ---- Standard examples -------------------------------------
\begin{exbox}[Example --- Linear Dependence and Independence]
\begin{enumerate}
    \item In $\R^2$: $\{(1,2),(-1,2)\}$ is l.i., since
          $\alpha(1,2)+\beta(-1,2)=\mathbf{0}$ gives $\alpha-\beta=0$ and
          $2\alpha+2\beta=0$, forcing $\alpha=\beta=0$.

    \item In $\F^n$: the standard basis vectors $\{e_1,\ldots,e_n\}$ are l.i.

    \item In $\F_n[x]$ (polynomials of degree $\le n$): $\{1,x,x^2,\ldots,x^n\}$
          is l.i., because a polynomial that is identically zero has all
          coefficients equal to zero.

    \item In $\F[x]$: $\{x^i\mid i\in\N_0\}$ is l.i.
\end{enumerate}
\end{exbox}

% ---- Method: row reduction ---------------------------------
\begin{methodbox}[Method --- Checking Linear Independence via Row Reduction]
\label{meth:li-row-reduction}
To test whether $v_1,\ldots,v_n\in\F^m$ are l.i.:
\begin{enumerate}
    \item Form the matrix $A=[v_1\mid\cdots\mid v_n]$ (vectors as \emph{columns}).
    \item Row-reduce $A$ to echelon form.
    \item The vectors are \textbf{l.i.\ iff} every column is a pivot column,
          equivalently $\rk(A)=n$.
          A zero row (or a non-pivot column) signals linear dependence.
\end{enumerate}
\end{methodbox}

% ---- Key properties ----------------------------------------
\begin{propbox}[Proposition --- Properties of Linear Independence]
\label{prop:li-properties}
Let $V$ be a vector space over $\F$.
\begin{enumerate}[label=(\roman*)]
    \item Any set containing $\mathbf{0}$ is l.d.
    \item If $S$ is l.d.\ and $S\subseteq T$, then $T$ is l.d.
    \item If $S\subseteq T$ and $T$ is l.i., then $S$ is l.i.
    \item If $S=\{v_1,\ldots,v_n\}$ is l.d.\ with $n\ge 2$, then some $v_i$ is
          a linear combination of the others.
    \item $S$ is l.i.\ if and only if every \emph{finite} subset of $S$ is l.i.
\end{enumerate}
\end{propbox}

\begin{proofbox}[Proof --- Properties of Linear Independence]
\begin{enumerate}[label=(\roman*)]
    \item $1\cdot\mathbf{0}=\mathbf{0}$ with coefficient $1\ne 0$.
    \item A non-trivial relation among vectors in $S$ is also a non-trivial
          relation among vectors in $T$.
    \item (Contrapositive of (ii).)
    \item If $\sum_{i=1}^n \alpha_i v_i=\mathbf{0}$ with some $\alpha_j\ne 0$,
          then $v_j = -\alpha_j^{-1}\sum_{i\ne j}\alpha_i v_i$.
    \item By definition, l.d.\ involves a \emph{finite} subset; so $S$ is l.d.\
          iff some finite subset is l.d.
\end{enumerate}
\hfill$\blacksquare$
\end{proofbox}

% ---- Two-vector criterion ----------------------------------
\begin{exbox}[Example --- Linear Independence of Two and Three Vectors]
\begin{itemize}
    \item $\{u,v\}$ is l.i.\ $\iff$ $u\ne\mathbf{0}$ and $v\ne\alpha u$ for
          every $\alpha\in\F$ (i.e.\ $v$ is not a scalar multiple of $u$).
    \item $\{u,v,w\}$ is l.i.\ $\iff$ $\{u,v\}$ is l.i.\ and
          $w\notin\spn\{u,v\}$.
\end{itemize}
\end{exbox}

% ---- Criterion for linear independence ---------------------
\begin{propbox}[Proposition --- Sequential Criterion for Linear Independence]
\label{prop:li-sequential}
A finite list $v_1,\ldots,v_r\in V$ is l.i.\ if
\begin{enumerate}[label=(\roman*)]
    \item $v_1\ne\mathbf{0}$, and
    \item $v_i\notin\spn\{v_1,\ldots,v_{i-1}\}$ for each $i=2,\ldots,r$.
\end{enumerate}
\end{propbox}

\begin{proofbox}[Proof --- Sequential Criterion for Linear Independence]
Suppose $\sum_{i=1}^r\alpha_i v_i=\mathbf{0}$.  Let $j$ be the largest index
with $\alpha_j\ne 0$.  If $j\ge 2$, then
$v_j = -\alpha_j^{-1}\sum_{i<j}\alpha_i v_i\in\spn\{v_1,\ldots,v_{j-1}\}$,
contradicting condition (ii).  If $j=1$ then $\alpha_1 v_1=\mathbf{0}$ with
$v_1\ne\mathbf{0}$, so $\alpha_1=0$.  Hence all $\alpha_i=0$.
\hfill$\blacksquare$
\end{proofbox}

% ---- Rank criterion for l.i. in F^m -----------------------
\begin{propbox}[Proposition --- Rank Criterion for Linear Independence]
\label{prop:rank-li}
Vectors $v_1,\ldots,v_n\in\F^m$ are l.i.\ if and only if
$\rk[v_1\mid\cdots\mid v_n]=n$.
\end{propbox}

\begin{proofbox}[Proof --- Rank Criterion for Linear Independence]
The vectors are l.i.\ iff $Ax=\mathbf{0}$ (with $A=[v_1\mid\cdots\mid v_n]$)
has only the trivial solution, iff the RREF of $A$ has no free variables, iff
every column of $A$ is a pivot column, iff $\rk(A)=n$.
\hfill$\blacksquare$
\end{proofbox}
